\section{Laplace-matched pressure self-duality}
The core kinetic pressure and circular Laplace pressure are
\[
  P_{\rm kin}=\frac12\rho_{\rm core}
  \lVert\mathbf v_{\!\boldsymbol{\circlearrowleft}}\rVert^2,
  \qquad
  P_{\rm Laplace}=\frac{F_{\rm swirl}^{\max}}{2\pi r_c^2}.
\]
If these agree, the inner and outer pressure moduli are equal at leading order.
The reciprocal matched action is then
\[
  A_\chi(\chi_R)=\chi_R+\frac{N_p}{\chi_R}.
\]
More generally,
\[
  A_\chi(\chi_R)=\chi_R+\mu\frac{N_p}{\chi_R},
  \qquad
  \mu=\frac{P_{\rm in}}{P_{\rm out}}.
\]
Stationarity gives
\[
  \chi_R=\sqrt{\mu N_p}.
\]
Thus for \(N_p=4\) and \(\mu=1\),
\[
  \chi_R=2,\qquad \lambda_\chi=4.
\]
This derives the self-dual radius action within the Laplace-matched reciprocal
pressure model.  A full primitive-equation derivation must still exclude
same-order nonreciprocal terms such as \(\log\chi_R\), \(\chi_R^2\), or
\(\chi_R^{-2}\), or show that they are higher-order corrections.
